\chapter{Analisi delle Tecnologie e Stato dell'Arte}
\label{chap:stato_arte}

% In questo capitolo, analizzi le tecnologie che hai considerato e le confronti con il tuo approccio.
La scelta di un'architettura per la gestione dei contenuti web richiede un'attenta analisi delle soluzioni esistenti, al fine di comprendere i compromessi di ciascuna e posizionare correttamente il proprio lavoro. Questo capitolo analizza il panorama tecnologico, dalle soluzioni più tradizionali a quelle più moderne, per poi approfondire il processo di valutazione e selezione degli strumenti specifici che ha portato all'architettura finale del progetto.

\section{Panoramica delle Soluzioni Esistenti}

\subsection{Sistemi di Content Management System (CMS) Tradizionali}
I CMS monolitici come WordPress, Joomla e Drupal rappresentano da anni la soluzione dominante per la creazione di siti web dinamici. Il loro punto di forza risiede in un'interfaccia di amministrazione completa e user-friendly, che astrae l'utente da ogni complessità tecnica, e in un vasto ecosistema di plugin e temi che estendono le funzionalità di base. 

Tuttavia, questa architettura "all-in-one" presenta svantaggi significativi che li rendono inadatti per il contesto di questa tesi: la dipendenza da uno stack server specifico (tipicamente LAMP/LEMP) e da un database relazionale li rende pesanti, complessi da manutenere e vulnerabili dal punto di vista della sicurezza. La necessità di aggiornamenti costanti sia del core che dei componenti terzi li trasforma in un bersaglio primario per attacchi informatici, richiedendo competenze specifiche per la loro manutenzione.

\subsection{Generatori di Siti Statici (SSG)}
In risposta alla complessità dei CMS tradizionali, sono emersi i Generatori di Siti Statici (SSG) come Jekyll, Hugo, Gatsby e Next.js (in modalità statica). Questi strumenti adottano un approccio fondamentalmente differente: separano la fase di costruzione del sito da quella di hosting. I contenuti, tipicamente scritti in formati semantici come Markdown o YAML, vengono processati staticamente per generare un sito composto unicamente da file HTML, CSS e JavaScript. 

I vantaggi sono notevoli in termini di performance, sicurezza (mancando un database o codice server-side dinamico, la superficie di attacco si riduce drasticamente) e semplicità di deployment. Il loro workflow è intrinsecamente basato su Git, favorendo pratiche di sviluppo moderne. Tuttavia, il loro principale limite, nel contesto di questa tesi, è l'assenza nativa di un'interfaccia di editing visuale online: la modifica dei contenuti richiede competenze tecniche (conoscenza di Git, formati di markup) e un processo manuale di modifica dei file sorgente e successivo commit, risultando poco accessibile per utenti non tecnici.

\subsection{Headless CMS}
L'approccio Headless CMS rappresenta un'evoluzione che tenta di unire i vantaggi dei due mondi precedenti. Sistemi come Contentful, Strapi, TinaCMS o Directus disaccoppiano completamente il backend di gestione dei contenuti (il "corpo") dal frontend di presentazione (la "testa"). I contenuti vengono gestiti tramite un'interfaccia web dedicata e resi disponibili tramite API, tipicamente REST o GraphQL, che il frontend consuma per renderizzare i contenuti.

Questo approccio offre massima flessibilità al frontend, permettendo di utilizzare qualsiasi tecnologia, e centralizza la gestione dei contenuti. Tuttavia, introduce nuove complessità: le soluzioni SaaS (Software as a Service) legano a un fornitore esterno con costi ricorrenti e dipendenze infrastrutturali, mentre quelle self-hosted possono richiedere una configurazione non banale (server, database, manutenzione). Inoltre, il flusso di editing è tipicamente disaccoppiato dalla visualizzazione finale, limitando l'esperienza WYSIWYG.

\section{Analisi e Scelta degli Strumenti}

La selezione degli strumenti per questo progetto è stata un processo iterativo, guidato dai requisiti specifici della traccia e confermato dall'esperienza pratica. Particolare attenzione è stata posta sulla scelta dell'editor WYSIWYG, componente cruciale per l'esperienza utente finale.

\subsection{La Scelta dell'Editor WYSIWYG}

Il percorso di selezione dell'editor è emblematico del processo di ricerca e sviluppo del progetto, dimostrando come requisiti specifici possano guidare verso scelte tecnologiche non ovvie.

\textbf{Valutazione iniziale: TinyMCE}
La prima scelta è ricaduta su TinyMCE, un editor JavaScript open-source estremamente potente, maturo e considerato uno standard industriale, utilizzato in progetti di grande rilevanza. La sua integrazione iniziale nel template Bootstrap "Nova" ha dimostrato la sua flessibilità nel rendere editabile qualsiasi elemento del DOM attraverso una configurazione minimale.

Tuttavia, sono emerse rapidamente delle complessità architetturali che lo rendevano non ottimale per questo specifico caso d'uso. TinyMCE è progettato primariamente per operare come un "campo di testo ricco" tradizionale e la sua integrazione in un flusso di salvataggio basato su file system e Git avrebbe richiesto una notevole quantità di codice custom:
\begin{itemize}
    \item \textbf{Identificazione del contenuto}: Sarebbe stato necessario sviluppare una logica complessa per identificare univocamente ogni area modificata e associarla a una porzione specifica del file HTML da salvare.
    \item \textbf{Pulizia dell'HTML}: L'output di TinyMCE, per quanto configurabile, tende a inserire classi e attributi specifici dell'editor (ad esempio per la gestione dello stile), che avrebbero dovuto essere "puliti" lato server prima di ogni salvataggio per non "inquinare" il codice sorgente del template.
    \item \textbf{Complessità lato client}: La gestione di più aree editabili indipendenti sulla stessa pagina avrebbe richiesto la configurazione e la gestione di multiple istanze dell'editor, aumentando significativamente la complessità del codice JavaScript lato client.
\end{itemize}

\textbf{La scelta finale: SimplyEdit}
Durante questa fase di valutazione, è emersa la libreria SimplyEdit. Questo strumento, sebbene meno noto di TinyMCE, si è rivelato più allineato alla filosofia del progetto per diverse ragioni fondamentali. 

A differenza di TinyMCE, SimplyEdit non si concepisce come un editor da "embeddare" in campi specifici, ma come uno strato applicativo che si sovrappone all'intera pagina. La sua filosofia "HTML-first", basata sull'aggiunta di semplici attributi \texttt{data-simply-*} direttamente nel codice HTML, ha permesso di definire le aree e le liste editabili in modo dichiarativo e con una quantità minima di JavaScript configurabile.

Il vantaggio decisivo è stata la sua architettura di storage intrinsecamente flessibile. SimplyEdit è progettato per essere agnostico rispetto al sistema di salvataggio sottostante e permette di sovrascrivere facilmente il suo storage layer di default. Questa caratteristica si è rivelata il gancio fondamentale che ha permesso di intercettare l'evento di salvataggio e di reindirizzarlo a un endpoint custom sul server Node.js, abilitando così l'integrazione diretta e pulita con il workflow di commit su Git, senza dover manipolare pesantemente l'output dell'editor.

\subsection{La Scelta della Piattaforma Backend}

Per il backend, la scelta è ricaduta su Node.js con il framework Express.js. Questa piattaforma si è rivelata ideale per il ruolo di "server leggero" richiesto dal progetto, offrendo diversi vantaggi chiave:

\begin{itemize}
    \item \textbf{Architettura a eventi non bloccante}: La natura asincrona di Node.js è perfetta per gestire operazioni di I/O come la scrittura di file sul disco, senza bloccare il thread principale e mantenendo alta la responsività del server.
    \item \textbf{Gestione efficiente di processi esterni}: La facilità con cui Node.js, tramite il modulo \texttt{child\_process}, può eseguire comandi di sistema è stata essenziale per l'integrazione trasparente con Git, permettendo di eseguire commit automatici in modo programmatico.
    \item \textbf{Ecosistema NPM}: L'accesso al vasto registro di pacchetti di NPM ha semplificato notevolmente l'implementazione di funzionalità avanzate come la gestione delle sessioni utente (\texttt{express-session}), il parsing HTML per lo scraping (\texttt{cheerio}), e il routing avanzato.
    \item \textbf{Isomorfismo JavaScript}: L'utilizzo di JavaScript sia lato client che lato server ha favorito una maggiore coerenza dello stack tecnologico e ha semplificato il processo di sviluppo.
\end{itemize}

\section{Posizionamento del Progetto}

Il sistema sviluppato in questa tesi occupa una posizione distintiva nel panorama delle soluzioni di content management, combinando strategicamente elementi di diverse categorie per creare un approccio ibrido che risponde in modo mirato a tutti i requisiti del progetto.

A differenza dei CMS tradizionali, la soluzione proposta è leggera, non richiede un database dedicato e minimizza la superficie di attacco. Contrariamente agli SSG puri, offre un'esperienza di editing visuale in tempo reale senza richiedere competenze tecniche agli editor dei contenuti. Rispetto ai headless CMS, è interamente self-hosted, non dipende da servizi esterni a pagamento e mantiene il contenuto strettamente accoppiato con la presentazione durante la fase di editing, garantendo una vera esperienza WYSIWYG.

La scelta di utilizzare Git non solo come sistema di versioning ma come strato di persistenza principale rappresenta un elemento particolarmente innovativo, che traccia un ponte ideale tra le best practices dello sviluppo software moderno e le esigenze pratiche della gestione contenuti. Questo approccio trasforma di fatto ogni modifica in un evento versionato e tracciabile, abilitando potenzialmente workflow avanzati di revisione, approvazione e deployment continuo.

Il progetto dimostra quindi come sia possibile creare un sistema di gestione contenuti dinamico, user-friendly e al tempo stesso tecnologicamente semplice, sicuro e mantenibile, colmando efficacemente il gap tra la rigidità dei siti statici e la complessità dei CMS tradizionali.